\documentclass[11pt,a4paper]{article}

\usepackage[margin=1in]{geometry}
\usepackage{times}
\usepackage{amsmath,amssymb}
\usepackage{booktabs}
\usepackage{longtable}
\usepackage{array}
\usepackage{xcolor}
\usepackage{graphicx}
\usepackage{tikz}
\usetikzlibrary{arrows.meta,positioning}
\usepackage{listings}
\usepackage[hidelinks]{hyperref}
\usepackage{caption}
\captionsetup{font=small,labelfont=bf}
\usepackage{enumitem}
\setlist{itemsep=2pt,topsep=3pt}

\definecolor{codebg}{RGB}{246,246,246}
\lstset{basicstyle=\ttfamily\footnotesize, backgroundcolor=\color{codebg},
        frame=single, rulecolor=\color{black!25}, breaklines=true,
        columns=fullflexible, xleftmargin=4pt, framexleftmargin=4pt}

\newcommand{\rr}[1]{\textcolor{red!70!black}{#1}}

\title{\textbf{Zero-Shot Open-Vocabulary Scene Graph Pipelines on\\
BIM-Derived RGB-D: A Controlled Study}\\[0.4em]
\large BIM-Graphs Technical Report}
\author{BIM-Graphs Project}
\date{\today}

\begin{document}
\maketitle

\begin{abstract}
\noindent
We evaluate whether open-vocabulary 3D scene graph pipelines can be applied
zero-shot to RGB-D flythroughs rendered from Building Information Models. Using
IFC as a source gives something no existing benchmark has: \emph{exact} ground
truth, since every pixel is traceable to a building element by construction. Our
first result was that ConceptGraphs performs very poorly on such data. Our
second, and more important, result is that this conclusion is \emph{not} safe
without a control: a matched run on Replica shows the same detector firing 25
objects per frame at confidence 0.49, against 1.93 at 0.26 on our renders, with
32\% of our frames yielding no detection at all. The deficit decomposes into two
independent causes. \textbf{(i)~Appearance}: IFC supplies one flat colour per
element, leaving 79\% of each frame featureless against Replica's 21\%; adding
procedural surface texture reduces empty frames from 31.7\% to 3.3\% and nearly
triples detections per frame. \textbf{(ii)~Content}: Replica's top detections are
pillow, sofa and coffee table, while ours are wall, door and window. A building
model is largely unfurnished architecture, and object-centric pipelines discard
walls and floors as background by construction --- so they search for precisely
what BIM omits and ignore what it states exactly. This is a structural
mismatch that no rendering fidelity can repair. We quantify both causes with a
three-condition appearance ladder, release the benchmark generator, and argue
that room- and topology-level metrics, not object recall, are the meaningful
measure for building-scale scene graphs.
\end{abstract}

\section{Motivation}

Open-vocabulary 3D scene graph methods are evaluated almost exclusively on
ScanNet, Replica and Matterport, whose ground truth is hand-annotated and
therefore imperfect and expensive. Evaluation consequently stops at object-level
metrics, and \emph{connectivity} --- which room opens onto which, the property a
robot actually navigates with --- is essentially never scored, because no
dataset states it reliably.

An IFC building model changes this. It provides exact geometry, exact element
identity, and (after the topology recovery described in our companion
literature review) exact room boundaries and connectivity. Rendering a
flythrough from it therefore yields posed RGB-D \emph{with} pixel-perfect
instance masks and a fully known scene graph. That makes controlled ablation
possible: pose noise, depth noise, appearance realism and furnishing density can
each be varied independently, which cannot be done on a real scan.

This report asks the prerequisite question: \emph{do existing zero-shot pipelines
work on such data at all, and if not, why not?}

\section{Benchmark construction}

\subsection{Ground truth generation}

For each IFC model we build the hybrid hierarchical scene graph (rooms,
elements, containment, door and passage connectivity), then plan a coverage
trajectory that walks the recovered navigation graph. Two properties of the
trajectory matter for fairness:

\begin{itemize}
\item It routes \emph{through doors} using the recovered topology. Straight
segments between room centres pass through walls, producing frames pressed
against masonry that would penalise any method for reasons unrelated to
perception.
\item Frames whose median depth falls below 0.8\,m are discarded, for the same
reason.
\end{itemize}

Each frame stores RGB, uint16 depth in millimetres, a per-pixel instance mask
keyed to the IFC \texttt{GlobalId}, the camera pose, and the room the camera
occupied. Rendering is offscreen via EGL at $640\times480$ with a $58^\circ$
vertical field of view, comparable to a RealSense D435.

\subsection{Appearance conditions}

Because we control the renderer, appearance realism becomes an experimental
variable rather than a confound. We define three conditions:

\begin{center}
\small
\begin{tabular}{@{}llp{8.2cm}@{}}
\toprule
\textbf{A} & \textsc{flat} & IFC diffuse colour only; one flat RGB per element. \\
\textbf{B} & \textsc{textured} & Procedural surface textures (plaster, wood,
tile, carpet, glass, metal) selected by IFC class and material name, with UVs in
world metres so pattern scale is physically consistent, plus oblique fill
lights. \\
\textbf{C} & \textsc{textured+clutter} & As B, plus procedural furniture placed
by room label. 98.5\% of rooms in the corpus carry a functional label, so a
\textsc{kitchen} can be furnished differently from a \textsc{bedroom}. Every
placed item is registered as a ground-truth instance. \\
\bottomrule
\end{tabular}
\end{center}

\subsection{Reference dataset}

As a control we use Replica \texttt{room0} in the Nice-SLAM RGB-D rendering ---
the dataset ConceptGraphs itself was developed against. It is run through the
identical detector, class list and configuration.

\section{Experimental setup}

\paragraph{Baselines.} ConceptGraphs (object-centric open-vocabulary,
YOLO-World + MobileSAM + CLIP ViT-H-14) and HOV-SG (hierarchical, with an
explicit room layer). Both are the published implementations, not
reimplementations.

\paragraph{Metrics.} We report detector-level statistics (fraction of frames
with no detection, detections per frame, mean confidence, distinct classes) and
scene-graph-level statistics against exact ground truth.

Two decisions make the scene-graph metrics meaningful:

\begin{itemize}
\item \textbf{Only observed ground truth is scored.} An element the camera never
saw cannot be recovered by any method; counting it as a miss would measure the
trajectory. Visibility is read from the rendered instance masks.
\item \textbf{Objects and architecture are scored separately.} ConceptGraphs
discards wall/floor/ceiling as background \emph{by design}, so scoring it against
a ground-truth set dominated by walls measures a category mismatch, not method
quality. We never merge the two buckets into one headline number.
\end{itemize}

We additionally report a \emph{hallucination rate}: the fraction of predicted
objects whose label is impossible in a building model (\texttt{person},
\texttt{dog}, \dots). This requires no matching and admits no ambiguity.

\section{Results}

\subsection{The control: is it the method or the renderer?}
\label{sec:control}

Table~\ref{tab:probe} is the central result of this report.

\begin{table}[h]
\centering
\small
\begin{tabular}{@{}lrrrrr@{}}
\toprule
\textbf{Source} & \textbf{Frames} & \textbf{Empty frames} &
\textbf{Det./frame} & \textbf{Confidence} & \textbf{Classes} \\
\midrule
A \textsc{flat} (model\_0)   & 120 & 31.7\% & 1.93 & 0.26 & 27 \\
B \textsc{textured}        &  60 &  3.3\% & 5.15 & 0.25 & 14 \\
\midrule
\textbf{Replica (reference)} & 120 & \textbf{0.0\%} & \textbf{25.02} &
\textbf{0.49} & \textbf{60} \\
\bottomrule
\end{tabular}
\caption{YOLO-World response under identical model, class list and thresholds.
The detector never fails on Replica and fires 13$\times$ more often there than
on flat BIM renders, at twice the confidence. A conclusion drawn from the flat
condition alone would have been a statement about our renderer.}
\label{tab:probe}
\end{table}

\subsection{Why: low-level image statistics}

\begin{table}[h]
\centering
\small
\begin{tabular}{@{}lrrr@{}}
\toprule
\textbf{Statistic} & \textbf{A \textsc{flat}} & \textbf{B \textsc{textured}} &
\textbf{Replica} \\
\midrule
Gradient energy        & 0.0031 & 0.0103 & 0.0208 \\
Local contrast         & 0.0177 & 0.0391 & 0.0763 \\
Saturation             & 0.070  & 0.150  & 0.161 \\
Distinct colours       & 35     & 91     & 306 \\
Intensity entropy      & 3.59   & 3.96   & 5.40 \\
\textbf{Featureless area} & \textbf{79\%} & \textbf{47\%} & \textbf{21\%} \\
\bottomrule
\end{tabular}
\caption{Frame statistics (means over sampled frames). ``Featureless area'' is
the fraction of $16\times16$ tiles with intensity standard deviation below
0.02. An IFC element carries a single diffuse colour, so a naive render is
overwhelmingly flat --- exactly the regime a detector trained on photographs has
never seen. Texturing recovers roughly half the deficit in gradient and
contrast, and matches Replica on saturation.}
\label{tab:stats}
\end{table}

\subsection{Scene-graph accuracy against exact ground truth}

\begin{table}[h]
\centering
\small
\begin{tabular}{@{}llrrrrrr@{}}
\toprule
\textbf{Cond.} & \textbf{Model} & \textbf{GT obj.} & \textbf{Pred.} &
\textbf{Recall} & \textbf{Prec.} & \textbf{Class acc.} & \textbf{Halluc.} \\
\midrule
A \textsc{flat}     & model\_0  &  87 &  40 & 0.333 & 0.725 & 0.241 & 0.075 \\
B \textsc{textured} & model\_0  &  87 &  45 & 0.333 & 0.644 & \textbf{0.483} & 0.067 \\
C \textsc{+clutter} & model\_0  & 161 &  88 & \textbf{0.466} & \textbf{0.852} & 0.440 & \textbf{0.011} \\
\midrule
A \textsc{flat}     & model\_17 & 104 &  59 & 0.406 & 0.729 & 0.209 & 0.000 \\
B \textsc{textured} & model\_17 & 104 &  82 & \textbf{0.510} & 0.646 & \textbf{0.528} & 0.000 \\
C \textsc{+clutter} & model\_17 & 169 & 118 & 0.509 & 0.729 & 0.419 & 0.000 \\
\bottomrule
\end{tabular}
\caption{ConceptGraphs against exact IFC ground truth, restricted to the object
bucket and to instances actually observed by the camera. Condition~C's ground
truth is larger because injected furniture is itself registered as ground truth,
so C-versus-B class accuracy is not like-for-like; recall, precision and
hallucination are.}
\label{tab:sg}
\end{table}

\paragraph{The two causes have distinct signatures.} Varying appearance and
content independently separates effects that a single ``more realistic
rendering'' change would have confounded:

\begin{center}
\small
\begin{tabular}{@{}lll@{}}
\toprule
& \textbf{Appearance (A$\to$B)} & \textbf{Content (B$\to$C)} \\
\midrule
Class accuracy & $\times$2 on both models & --- \\
Recall         & unchanged (model\_0)     & $+0.133$ (model\_0) \\
Precision      & $-0.08$ on both          & $+0.21$ (model\_0) \\
Hallucination  & barely moves             & $0.067 \to 0.011$ \\
\bottomrule
\end{tabular}
\end{center}

Texture fixes \emph{naming}; content fixes \emph{finding}. Neither substitutes
for the other.

\paragraph{Texture changes naming, not finding.} On model\_0 the detector-level
gain from texturing is large (1.93 to 5.15 detections per frame) yet the number
of matched ground-truth objects is identical at 29. The additional detections
are not new objects: they are more confident labels on regions already being
detected, plus additional false positives, which is why precision falls. Read
alone, the detector statistics would imply a recall improvement that did not
occur.

\paragraph{Clutter helps only where content was missing.} Adding furniture
raises recall by $0.133$ on model\_0 but by $0.001$ on model\_17. This is what
the content hypothesis predicts rather than a weakness of it: model\_17 is an
office whose rooms already contain furniture, so there was little deficit to
correct, whereas model\_0 is an open-corridor residential block whose coverage
path spends much of its time on an unfurnished exterior deck. A uniform
improvement across both buildings would have been the more suspicious result.

\paragraph{Hallucination tracks content, not appearance.} The impossible-label
rate on model\_0 falls $0.075 \to 0.067 \to 0.011$. The \texttt{person}
detections that survived filtering and merging in conditions A and B were not a
quirk of the detector: they are what an object-seeking model does when aimed at
empty architecture. Give it real objects and the phantoms largely stop.

\paragraph{Texture changes naming, not finding.} On model\_0 the detector-level
gain is large (1.93 to 5.15 detections per frame) yet the number of matched
ground-truth objects is identical at 29. The additional detections are not new
objects: they are more confident labels on regions already being detected, plus
additional false positives, which is why precision falls. Read alone, the
detector statistics would imply a recall improvement that did not occur.

\paragraph{Building type dominates appearance.} The gap between the two
buildings under identical rendering (recall 0.510 vs 0.333 in condition~B)
exceeds the gap between appearance conditions within either building. A single
corpus-wide average would therefore be misleading, and we report per-model
figures throughout.

\subsection{HOV-SG}

HOV-SG is the more meaningful comparison because it models rooms explicitly.
Its run is in progress at the time of writing: feature extraction over 51
keyframes completes in 4m43s, but sequential 3D mask merging requires
approximately 165\,s per iteration over 50 iterations ($\approx$2\,h15). We note
this because an earlier attempt was terminated by our own 50-minute wall clock
and could easily have been misread as a failure of the method.

One configuration detail is worth recording: \texttt{create\_graph.py} skips
hierarchical graph construction entirely for datasets named \texttt{replica} or
\texttt{scannet}, on the assumption that those are single-room scans. Because
our renders adopt the Replica on-disk layout, registering them under that name
would have produced a complete-looking run with the room layer --- the only
reason to use HOV-SG --- silently absent.

\subsection{The second cause: content, not appearance}
\label{sec:content}

Texture alone does not close the gap, and the class histograms show why:

\begin{center}
\small
\begin{tabular}{@{}lp{9.4cm}@{}}
\toprule
\textbf{Replica} & pillow (690), sofa chair (249), coffee table (219),
lamp (186), book (169), couch (147), shelf (142), potted plant (116) \\
\addlinespace
\textbf{BIM (B)} & closet door (107), closet wall (67), window (63),
picture (28), bed (14), person (11) \\
\bottomrule
\end{tabular}
\end{center}

Replica rooms are \emph{furnished}; IFC buildings are \emph{architecture}. Across
four storeys, model\_0 contains 89 furniture instances in total, while a single
Replica room presents hundreds of detectable objects. Since ConceptGraphs
explicitly discards wall, floor and ceiling as background, it is searching for
the category BIM does not model while ignoring the category BIM specifies
exactly.

\paragraph{Hallucination.} The label \texttt{person} appears in every BIM
condition. There are no people in an IFC model. In the flat condition three
\texttt{person} objects survived the three-detection filter \emph{and} the
merging stage, meaning the pipeline built stable 3D objects out of nothing. The
rate falls with texture but does not vanish.

\section{Discussion}

\paragraph{The negative result was not safe.} Our first measurement --- six
recovered objects for a four-storey building --- looked like decisive failure. It
was an artefact of frame stride: at 1.75\,m spacing the median object was seen
in exactly one frame, so multi-view association never fired. Dense sampling took
the count from 6 to 40 and the median from 1 to 6 detections per object. A
second apparent failure, near-total detector silence, proved to be flat
rendering. Both would have been reported as method limitations.

\paragraph{What the field should take from this.} Applying object-centric
open-vocabulary pipelines to BIM is a category error, independent of rendering.
The meaningful comparison is with methods that model \emph{rooms} --- HOV-SG and
its successors --- and the meaningful metric is room segmentation and
connectivity, not object recall. Object recall on an unfurnished building
measures the dataset.

\paragraph{Implication for our own pipeline.} This supports scoping
vision-language models to what IFC lacks --- affordances, and functional labels
for the 344 geometrically recovered rooms that carry no name --- rather than to
rediscovering objects that were never in the file.

\section{Threats to validity}

\begin{enumerate}
\item \textbf{Synthetic appearance remains unmatched.} Condition B closes about
half the statistical gap to Replica, not all of it. Absolute scores on BIM
renders should not be read as scores on real buildings.
\item \textbf{Idealised sensing cuts the other way.} Our poses are exact and
depth is noiseless, which \emph{flatters} the methods relative to a real robot.
Appearance and sensing biases therefore act in opposite directions and must not
be conflated.
\item \textbf{Single-detector control.} The Replica control isolates the
detector, which is where the failure begins, but the association and fusion
stages are not separately controlled.
\item \textbf{Procedural clutter is not photoreal.} Condition C tests the effect
of \emph{content presence}, not of realistic objects. It should be read as a
lower bound on what furnishing contributes.
\item \textbf{Scale.} Results reported here cover a subset of the corpus; the
generator runs over all 25 models.
\end{enumerate}

\section{Reproducibility notes}

Both baselines are the published implementations. Getting them to run required
resolving version drift that we record because it is likely to affect anyone
repeating this work.

\paragraph{ConceptGraphs.} Eight blockers: absent \texttt{nvcc}; CUDA~12.1
rejecting \texttt{gcc}$>$12; pip build isolation hiding \texttt{torch} from the
\texttt{chamferdist} build; conda CUDA headers living under
\texttt{targets/x86\_64-linux/include}; \texttt{supervision} removing
\texttt{ColorPalette.default()} (and being silently upgraded again by the
GroundingDINO install); \texttt{transformers} too new for RAM; missing
\texttt{fairscale}. Two defects were in the repository itself: a
\texttt{defaults:\ -\ base} entry naming a file that does not exist, and a
default \texttt{match\_method} the code does not implement.

\paragraph{HOV-SG.} numpy~2 breaking the compiled \texttt{faiss};
\texttt{scikit-fmm} requiring a C++ compiler; a hardcoded interactive
matplotlib backend.

\paragraph{Modifications to third-party code.} Four, all recorded in-line and
none algorithmic: skip frames where YOLO-World returns no boxes (ultralytics
SAM raises on an empty bbox tensor); skip the corresponding absent detection
files during mapping; switch HOV-SG's matplotlib backend to \texttt{Agg} for
headless operation; add a dataset branch so our Replica-layout data reaches
HOV-SG's hierarchical graph construction, which is otherwise skipped for
\texttt{replica} and \texttt{scannet}.

\section{Conclusions}

\begin{enumerate}
\item Zero-shot open-vocabulary pipelines do perform poorly on BIM-derived
RGB-D, but the naive measurement overstates it by a wide margin. Controls are
mandatory: without the Replica reference we would have attributed a rendering
artefact to the method.
\item \textbf{Appearance and content are separable causes with distinct
signatures.} Texture doubles class accuracy and leaves recall untouched; content
raises recall and precision and reduces hallucination by a factor of six.
Neither substitutes for the other, and a single ``more realistic rendering''
intervention would have confounded them.
\item Object-centric pipelines remain structurally mismatched to BIM: they
discard walls and floors as background while seeking the furniture a building
model does not contain. Injected clutter helps only where content was genuinely
absent.
\item Building-scale scene graph evaluation should be room- and
topology-centric. Exact IFC ground truth makes that measurable for the first
time.
\end{enumerate}

\paragraph{Errors caught during this study.} We record these because each
initially looked like a result. (i)~A frame stride of five left the median object
visible in one frame, suppressing multi-view association and yielding six
objects for a four-storey building. (ii)~Flat rendering silenced the detector on
a third of frames. (iii)~Detections per frame nearly tripled with texture while
matched objects did not change at all. (iv)~Our own evaluator scored correctly
identified clutter as class errors, making condition~C appear to regress from
0.483 to 0.187 when the true figure is 0.440. (v)~A 50-minute timeout killed
HOV-SG mid-merge. Five apparent findings, all artefacts of setup or measurement.

\end{document}
